\documentclass[11pt,a4paper]{article}
\usepackage{amsmath,amssymb,amsthm}
\usepackage{geometry}
\usepackage{hyperref}
\usepackage{booktabs}
\usepackage{array}
\geometry{margin=2.5cm}

\newtheorem{theorem}{Theorem}[section]
\newtheorem{lemma}[theorem]{Lemma}
\newtheorem{corollary}[theorem]{Corollary}
\newtheorem{definition}[theorem]{Definition}
\newtheorem{remark}[theorem]{Remark}
\newtheorem{observation}[theorem]{Observation}

\title{The Cascade Series --- Part I\\[0.5em]
\large Scale Variance from Orthogonality:\\
How the Unit Ball Generates $10^{120}$ Orders of Magnitude}
\author{[Author]}
\date{March 2026}

\begin{document}
\maketitle

\begin{abstract}
We prove that the unit ball in $\mathbb{R}^d$, equipped only with the slicing recurrence
$V_{d+1}/V_d = \sqrt{\pi} \cdot R(d)$, produces two distinguished geometric scales separated by
120 orders of magnitude. The mechanism is elementary: the recurrence contains a
unique dimension-independent constant, $\sqrt{\pi} = \Gamma(1/2)$, forced by orthogonality. The
sphere-area decay rate $p(d)$ has a natural zero at $d_0 \approx 2\pi$, the sphere-area maximum.
This zero-crossing condition forces a first threshold value $c_1 = \tfrac{1}{2}\ln\pi$ as the value
of the $d$-dependent part of $p$ at the zero. The exp-conjugacy of the recurrence's
two canonical forms then forces a second threshold value $c_2 = e^{c_1} = \sqrt{\pi}$. The
pair $\{c_1, c_2\}$ is the unique pair generated from $\{0\}$ by operations internal to the
recurrence, and no further element is reachable; the orbit terminates.

These two thresholds select two distinguished dimensions $d_1 = 19$ and $d_2 = 217$.
The sphere areas at these dimensions define two geometric scales: $\Omega_{19} \approx 0.52$
(an $O(1)$ number) and $\Omega_{217} \approx 10^{-120}$ (an exponentially small number). The
$\sqrt{\pi}$-dependent formula for 120 is: $\pi e^{2\sqrt{\pi}}(2\sqrt{\pi}-1)/\ln 10 \approx 120.3$. The cascade invariant
$I := \Omega_{19} \times \Omega_{217} \times 198/217 = 1.0996\times 10^{-120}$ is the natural multiplicative combination
of both threshold outputs; by Corollary~8.3 it is a function of the primitive pair $(c_1, c_2)$
alone, with zero free parameters. Its uniqueness is established (Theorem~9.2): under
one geometric principle---(A2) that the invariant uses the full primitive pair---five
constraints from the cascade's geometry admit exactly one scalar. The constraint
that content is measured by sphere areas with exponent 1 is not an independent
assumption: it follows from the recurrence's first-order multiplicative structure and
the fact that volumes are derived quantities ($V_d = \Omega_{d-1}/d$).

Every step from orthogonality to the threshold pair $(d_1, d_2)$ is a theorem about
the Gamma function. No physics enters the derivation. The numerical agreement of
$I$ with the observed vacuum-energy ratio $\rho_\Lambda/M^4_{\mathrm{Pl}} \approx 1.1\times 10^{-120}$ (to 0.04\%) is noted
as an observation, not explained within this paper.

Section~4 proves that each step of the slicing recurrence sends one direction's
effective radius to $R_{\mathrm{eff}} = 1/\sqrt{d+3}$, derived exactly from the Beta function. The
boundary-to-volume ratio $\Omega_{d-1}/V_d = d$ proves that sphere areas---not volumes---are
the cascade's primary objects. Corollary~4.12 establishes that sphere area is the only
independent cascade quantity at each level: every other cascade object is derived
from sphere areas, so any physical identification mapping cascade content to energy
must do so with a universal constant, with no level-type dependence available.
\end{abstract}

\tableofcontents
\newpage

\section{The Question}

How much scale variance can pure mathematics produce from nothing?

The unit ball $B^d = \{x \in \mathbb{R}^d : |x| \leq 1\}$ is the simplest geometric object that exists in
every dimension. Its volume is $V_d = \pi^{d/2}/\Gamma(d/2+1)$, and the surface area of its boundary
sphere $S^{d-1}$ is $\Omega_{d-1} = 2\pi^{d/2}/\Gamma(d/2)$. Both tend to zero as $d\to\infty$.

The sphere area $\Omega_d$ varies over an enormous range as $d$ varies. At $d=7$, $\Omega_7 \approx 33$
(near its maximum). At $d=217$, $\Omega_{217} \approx 10^{-120}$. The ratio is $\Omega_7/\Omega_{217} \approx 10^{121}$. The unit
ball, with no free parameters and no input beyond the definition of orthogonal dimensions,
generates 121 orders of magnitude of scale variance.

This paper asks: is there a natural way to extract two distinguished scales from this
variance? The answer is yes. The slicing recurrence that relates $V_{d+1}$ to $V_d$ contains exactly
one dimension-independent constant ($\sqrt{\pi}$). The sphere-area decay rate has a natural
zero (the sphere-area maximum), and this zero canonically generates two threshold values via
operations internal to the recurrence. These thresholds select exactly two dimensions (19
and 217) and their sphere areas. The two scales are not arbitrary choices from the $\Omega_d$
profile. They are the unique outputs of the recurrence's internal structure, rooted in its
natural zero.

\section{Rates as Fundamental Objects}

The scale variance of $\Omega_d$ is most naturally studied through its rate of change. As $d\to\infty$,
volumes and surface areas are identically zero: all measure concentrates on a boundary
layer of thickness $\sim 1/\sqrt{d}$, and the boundary itself has zero area. The only well-defined
quantities are rates.

\begin{definition}[Sphere-area decay rate]
$p(d) := -\partial_d\ln\Omega_d = \frac{1}{2}\psi\!\left(\frac{d+1}{2}\right) - \frac{1}{2}\ln\pi$,
where $\psi = \Gamma'/\Gamma$ is the digamma function.
\end{definition}

For large $d$, using $\psi(x) \approx \ln x$: $p(d) \approx \frac{1}{2}\ln(d/2\pi)$. This rate increases strictly and
continuously on $(0,\infty)$, from $-\infty$ to $+\infty$, with a unique zero at $d_0$ satisfying $p(d_0) = 0$
(the sphere-area maximum, $d_0 \approx 2\pi$).

\begin{remark}
The rate $p(d)$ is not the logarithm of some more fundamental quantity. At
$d=\infty$, where every $\Omega_d$ is zero, the rate of approach to zero is all the geometry there is. A
constant that enters the recurrence as a rate is already in the space where $p(d)$ lives. This
will be critical in Section~6.
\end{remark}

\begin{remark}[The three zeros]
The cascade has a single zero, appearing in three representations. (i) Geometric zero: $V_\infty = \Omega_\infty = 0$. The infinite-dimensional unit ball has
no volume, no surface area, no interior. This is the cascade's starting point---geometric
nothing. (ii) Rate-space zero: $p(d_0) = 0$. The dimension where growth transitions to
decay. Its uniqueness is a theorem ($p$ is strictly monotone). (iii) Orbit seed: $\{0\} \in \mathbb{R}$.
The starting point of the orbit construction in Section~6.

These three zeros are not the same object: the geometric zero occurs at $d=\infty$; the
rate zero occurs at $d_0 \approx 2\pi$; and the orbit seed is an algebraic value. They are, however,
coherently related. The cascade's geometric origin is zero, making zero the natural value
to encounter as the cascade's structural feature in rate-space. The orbit seed inherits its
justification from this structure: it is the value $p$ takes at the cascade's unique structural
zero, read off from the decay rate at $d_0$.

The practical consequence is that the orbit seed $\{0\}$ requires no external justification and
constitutes no additional input. It is read directly from the construction. Any alternative
seed would be an imported value not present in the cascade.
\end{remark}

\section{Orthogonal Compression}

\subsection{The slicing integral}

Cut $B^{d+1}$ perpendicular to one axis. Each cross-section at height $x$ is a $d$-ball of radius
$\sqrt{1-x^2}$:
\begin{equation}
V_{d+1} = V_d \int_{-1}^{1} (1-x^2)^{d/2}\,dx.
\end{equation}

\subsection{The quarter turn}

Substitute $x = \cos\theta$:
\begin{equation}
V_{d+1} = 2V_d \int_0^{\pi/2} \sin^d\theta\,d\theta.
\end{equation}
The range $[0,\pi/2]$ is one quarter turn: the rotation from the axis to the equator.

\begin{theorem}[Orthogonal compression constant]
$\Gamma(1/2) = \sqrt{\pi}$ is the unique dimension-independent constant produced by the quarter-turn integral. It is forced by orthogonality alone. Moreover, it appears with coefficient exactly 1 in each canonical form of the recurrence.
\end{theorem}

\begin{proof}
The quarter-turn integral evaluates via the Beta function:
\[
\int_0^{\pi/2}\sin^d\theta\,d\theta = \tfrac{1}{2}B\!\left(\tfrac{1}{2},\tfrac{d+1}{2}\right) = \frac{\Gamma(1/2)\,\Gamma((d+1)/2)}{2\,\Gamma((d+2)/2)}.
\]
The first argument of $B$ is $\frac{1}{2}$, forced by orthogonality: the angle between axis and equator
is $\pi/2$; the substitution $x = \cos\theta$ maps $[0,\pi/2]$ to $[0,1]$, producing the half-power $t^{-1/2}$
in $B(1/2,\cdot)$; this gives $\Gamma(1/2) = \sqrt{\pi}$. The chain is: dimension $\to$ orthogonal $\to$ $\pi/2$ $\to$
$B(1/2,\cdot)$ $\to$ $\Gamma(1/2) = \sqrt{\pi}$. Every arrow is forced.

For the coefficient: the linear-form recurrence is $V_{d+1}/V_d = \sqrt{\pi}\cdot R(d)$ where $R(d) = \Gamma((d+1)/2)/\Gamma((d+2)/2)$. The factor $\sqrt{\pi}$ appears once, with no additional numeric coefficient. Taking logarithms: $\Delta\ln V_d = \frac{1}{2}\ln\pi + \ln R(d)$. The constant term $\frac{1}{2}\ln\pi$ has coefficient 1, and the $d$-dependent part $\ln R(d)$ contains no $\pi$-dependent constants by Theorem~5.1 below.
\end{proof}

\subsection{The full slicing recurrence}

Combining both hemispheres:
\begin{equation}
V_{d+1} = V_d\cdot\sqrt{\pi}\cdot R(d), \qquad R(d) := \frac{\Gamma((d+1)/2)}{\Gamma((d+2)/2)}.
\end{equation}
For large $d$, $R(d) \approx \sqrt{2/d}$, giving $V_{d+1}/V_d \approx \sqrt{2\pi/d}$.

\section{Asymptotic Compactification}

This section has a single goal: to establish that sphere areas are the unique independent cascade quantity at each level (Corollary~4.12). All intermediate results serve that end. Theorem~4.2 derives the compactification radius $R_\mathrm{eff} = 1/\sqrt{d+3}$ from the Beta function; Theorem~4.4 shows it coincides asymptotically with the Gaussian width; Theorem~4.7 proves boundary dominance $\Omega_{d-1}/V_d = d$; Remark~4.11 uses that identity to show volumes are derived quantities. Corollary~4.12 assembles these into the statement that every cascade object is a function of sphere areas alone, and that any physical identification mapping cascade content to energy must use a universal constant with no level-type dependence. Readers who accept this conclusion may proceed directly to Section~5.

\subsection{Concentration of the integrand}

\begin{lemma}[Gaussian concentration]
The integrand $f_d(x) = (1-x^2)^{d/2}$ satisfies:
(a) $f_d(x) = \exp(-dx^2/2 + O(dx^4))$, approximating a Gaussian with standard deviation
$\sigma(d) = 1/\sqrt{d}$; (b) the half-maximum width is $\sqrt{2\ln 2}/\sqrt{d}$; (c) as $d\to\infty$, $f_d\to\delta(x)$.
\end{lemma}

\begin{proof}
Expanding: $\ln(1-x^2) = -x^2 - x^4/2 - x^6/3 - \cdots$ for $|x|<1$. Therefore $(d/2)\ln(1-x^2) = -dx^2/2 + O(dx^4)$. The half-maximum occurs at $x_{1/2} = \sqrt{1-2^{-2/d}} \approx \sqrt{2\ln 2}/\sqrt{d}$.
As $d\to\infty$, $x_{1/2}\to 0$ and $\int f_d\,dx \approx \sqrt{2\pi/d}\to 0$, confirming convergence to $\delta(x)$.
\end{proof}

\subsection{The compactification radius}

\begin{theorem}[Effective radius]
The effective radius of the compactified direction at step $d$ is $R_{\mathrm{eff}}(d) = 1/\sqrt{d+3}$.
\end{theorem}

\begin{proof}
The effective radius is the RMS extent: $R_{\mathrm{eff}}^2 = \langle x^2\rangle = \int x^2(1-x^2)^{d/2}dx\Big/\int(1-x^2)^{d/2}dx$. Evaluating via the Beta function with $t = x^2$: the numerator is $B(3/2, d/2+1)$ and the denominator is $B(1/2, d/2+1)$. Using $\Gamma(x+1) = x\Gamma(x)$: $\langle x^2\rangle = 1/(d+3)$. Therefore $R_{\mathrm{eff}} = 1/\sqrt{d+3}\to 0$ as $d\to\infty$.
\end{proof}

\begin{center}
\begin{tabular}{cccl}
\toprule
$d$ & $R_{\mathrm{eff}}(d)$ & $N(d)$ & Layer\\
\midrule
4 & 0.37796 & 1.17810 & Observer\\
7 & 0.31623 & 0.91429 & $d_0$\\
12 & 0.25820 & 0.70870 & $d=12$\\
19 & 0.21320 & 0.56755 & $d_1$\\
217 & 0.06742 & 0.16997 & $d_2$\\
\bottomrule
\end{tabular}
\end{center}

\begin{remark}
The compactification is smooth. At finite $d$, the weight $(1-x^2)^{d/2}$ has
support on all of $(-1,1)$. The direction is not removed; its contribution is suppressed to
$O(1/\sqrt{d})$. There is no boundary, no horizon, no sharp cutoff---only exponential suppression.
The word ``asymptotic'' is essential: the compactification limit is approached but never
reached at any finite step.
\end{remark}

\subsection{The effective-radius and Gaussian-width coincidence}

\begin{theorem}[Scale coincidence]
The cascade's two natural length scales---the compactification radius $R_{\mathrm{eff}}(d) = 1/\sqrt{d+3}$ and the Gaussian width $\sigma(d) = 1/\sqrt{d}$ of the slicing integrand---satisfy
\[
\frac{R_{\mathrm{eff}}(d)}{\sigma(d)} = \sqrt{\frac{d}{d+3}} \to 1 \quad\text{as } d\to\infty,
\]
with correction $O(1/d)$: $R_{\mathrm{eff}}/\sigma = 1 - \frac{3}{2d} + O(d^{-2})$.
\end{theorem}

\begin{proof}
Direct computation from Lemma~4.1 and Theorem~4.2.
\end{proof}

\begin{corollary}[Second universal cascade constant]
The fraction of the slicing integrand's content within the compactification radius satisfies
\[
F(d) = \frac{\int_{|x|<R_{\mathrm{eff}}(d)}(1-x^2)^{d/2}dx}{\int_{-1}^{1}(1-x^2)^{d/2}dx} \to \operatorname{erf}\!\left(\tfrac{1}{\sqrt{2}}\right) = 0.68269\ldots \quad\text{as }d\to\infty,
\]
a second universal constant of the cascade alongside $\sqrt{\pi}$, determined entirely by the coincidence $R_{\mathrm{eff}}/\sigma\to 1$.
\end{corollary}

\begin{proof}
In the Gaussian limit (Lemma~4.1), the integrand is $e^{-dx^2/2}$ with width $\sigma = 1/\sqrt{d}$.
The fraction within $R_{\mathrm{eff}} = 1/\sqrt{d+3}$ is $\operatorname{erf}(R_{\mathrm{eff}}/(\sigma\sqrt{2}))$. As $d\to\infty$, $R_{\mathrm{eff}}/(\sigma\sqrt{2})\to 1/\sqrt{2}$, giving $\operatorname{erf}(1/\sqrt{2})$.
\end{proof}

\begin{remark}[The cascade's effective geometric origin]
After $d_2 = 217$, sphere areas satisfy $\Omega_d < 10^{-120}$---geometrically indistinguishable from zero. The cascade's second threshold $d_2 = 217$ and its formal starting point $d=\infty$ are therefore the same geometric point: the cascade begins at geometric nothing, and nothing remains after $d_2$. The limit $d\to\infty$ in Corollary~4.5 is consequently exact, not merely asymptotic, when applied to the cascade tower: there is no geometric content beyond $d_2 = 217$ to distinguish the two points. The constant $\operatorname{erf}(1/\sqrt{2})$ is a property of the cascade's origin, evaluated at its effective terminus.
\end{remark}

\subsection{Why sphere areas are primary}

\begin{theorem}[Boundary dominance]
$\Omega_{d-1}/V_d = d$ for all $d\geq 1$.
\end{theorem}

\begin{proof}
$\Omega_{d-1} = 2\pi^{d/2}/\Gamma(d/2)$, $V_d = \pi^{d/2}/\Gamma(d/2+1) = \pi^{d/2}/[(d/2)\Gamma(d/2)]$. Ratio: $2\times(d/2) = d$.
\end{proof}

\begin{corollary}[Sphere areas as boundary measures]
The sphere area $\Omega_d$ is the natural measure of the cascade's content at dimension $d$, because $\Omega_{d-1}/V_d = d$: at high $d$, the boundary dominates the volume. By concentration of measure, the unit ball's volume concentrates in a shell of thickness $\sim 1/\sqrt{d}$ near $S^{d-1}$. This is the reason the cascade's thresholds (Section~6) and invariant (Section~9) are defined in terms of sphere areas rather than volumes.
\end{corollary}

\subsection{The cascade as a compactification sequence}

\begin{theorem}[Compactification sequence]
The cascade $V_\infty\to V_{d_2}\to V_{d_1}\to V_{d_0}\to V_4$ has at each step $d+1\to d$: (a) effective radius $1/\sqrt{d+3}$ (Theorem~4.2); (b) volume multiplied by $N(d)<1$ for $d>d_0$; (c) boundary-to-volume ratio $\Omega/V = d$ (Theorem~4.7); (d) the integrand $(1-x^2)^{d/2}$ is smooth and positive on $(-1,1)$ for all finite $d$.
\end{theorem}

\subsection{Layered boundary dominance}

\begin{corollary}[Layered boundary dominance]
The cascade's boundary-to-volume ratio $\Omega/V = d$ (Theorem~4.7), combined with the compactification radius $R_{\mathrm{eff}} = 1/\sqrt{d+3}$ (Theorem~4.2), implies: (a) at each step, the boundary $S^{d-1}$ carries a fraction $d/(d+1)$ of the content; (b) the boundary capacity $\Omega_d$ is calculable from $\pi$ alone; (c) the cascade stacks these boundary-dominant steps: the $d_2$-sphere's boundary encloses the $d_1$-sphere's boundary, which encloses the $d_0$-sphere's boundary.
\end{corollary}

\begin{remark}[Volumes as derived quantities]
Since $V_d = \Omega_{d-1}/d$ is an identity (not a physical relationship), every quantity in this paper---the decay rate $p(d)$, the threshold conditions, the cascade invariant---is expressible purely in terms of sphere areas. Volumes have no independent content; they are boundary areas divided by the dimension. The slicing recurrence, written as
\[
\Omega_d = \Omega_{d-1}\cdot\frac{d+1}{d}\cdot\sqrt{\pi}\cdot R(d),
\]
is a first-order multiplicative map between sphere areas. The integral $\int(1-x^2)^{d/2}dx$ is the kernel of this map: it computes the sphere-area ratio $\Omega_d/\Omega_{d-1}$ (up to the dimensional prefactor), not an independent volume.
\end{remark}

\begin{corollary}[Sphere areas are the unique independent cascade quantities]
Every cascade quantity---$V_d$, $N(d)$, $R(d)$, $R_{\mathrm{eff}}(d)$---is derived from sphere areas:
\[
V_d = \frac{\Omega_{d-1}}{d}, \quad N(d) = \frac{\Omega_d}{\Omega_{d-1}}, \quad R(d) = \frac{N(d)}{\sqrt{\pi}}, \quad R_{\mathrm{eff}}(d) = \sqrt{\langle x^2\rangle} = \frac{1}{\sqrt{d+3}}.
\]
No independent cascade object exists at level $d$ other than $\Omega_{d-1}$. Any physical identification that maps cascade content to energy must therefore assign energy $C\cdot\Omega_{d-1}$ at each level, where $C$ is a universal constant: no independent geometric quantity exists to make $C$ depend on $d$ or layer type.
\end{corollary}

\begin{proof}
$V_d = \Omega_{d-1}/d$ is the identity $\Omega_{d-1}/V_d = d$ (Theorem~4.7) solved for $V_d$. $N(d) = V_{d+1}/V_d = \Omega_d/\Omega_{d-1}$ after applying the identity at both levels. $R(d) = N(d)/\sqrt{\pi}$ by the recurrence. $R_{\mathrm{eff}}(d) = 1/\sqrt{d+3}$ from Theorem~4.2. Every cascade quantity is therefore a function of sphere areas alone.

Since $\Omega_{d-1}$ is the content of level $d$---not a proxy for it---and any physical identification maps content directly to energy, the energy at level $d$ is $C\cdot\Omega_{d-1}$ by direct substitution. The recurrence structure confirms this: the slicing recurrence is a first-order multiplicative map outputting exactly one factor of $\Omega_d$ per step (via $N(d) = \Omega_d/\Omega_{d-1}$), so the natural multiplicity of each level's contribution is 1. The constant $C$ is universal across all levels and layer types.
\end{proof}

\section{The Irreducibility of $\sqrt{\pi}$}

\begin{theorem}[Gamma-ratio exclusion]
Define $\delta_{\mathrm{exact}}(d) := \ln(\Omega_{d+1}/\Omega_d) - \ln(\Omega_d/\Omega_{d-1})$. Then $\delta_{\mathrm{exact}}(d) = \ln(d/(d+1)) < 0$ for all $d\geq 1$.
\end{theorem}

\begin{proof}
Using $\Omega_d = (d+1)V_{d+1}$ and the recurrence: the Gamma functions cancel pairwise under $\Gamma(x+1) = x\Gamma(x)$, leaving $\ln(d/(d+1)) < 0$.
\end{proof}

\begin{corollary}[Irreducibility]
The second difference of $\ln\Omega_d$ reduces to a rational expression. All $\pi$-dependent and $\Gamma$-dependent factors cancel exactly. $\sqrt{\pi}$ is the irreducible residue: the only nontrivial constant the sphere-area recurrence contains.
\end{corollary}

\section{The Natural Zero Generates the Threshold Pair}

This is the central section of the paper.

\subsection{Decomposition of the decay rate}

The decay rate decomposes canonically as:
\[
p(d) = -\tfrac{1}{2}\ln\pi + \tfrac{1}{2}\psi\!\left(\tfrac{d+1}{2}\right).
\]
The constant part is $-\frac{1}{2}\ln\pi$, arising directly from the $\sqrt{\pi}$ in the recurrence. The $d$-dependent part $q(d) = \frac{1}{2}\psi((d+1)/2)$ is strictly increasing from $-\infty$ to $+\infty$.

\subsection{The natural zero}

\begin{definition}[Natural zero]
The natural zero $d_0$ of the recurrence is the unique solution to $p(d_0) = 0$: $q(d_0) = \frac{1}{2}\ln\pi$. In the Stirling approximation, $d_0 = 2\pi\approx 6.28$, integer $d_0 = 7$. At $d_0$, the sphere-area sequence transitions from growth to decay: $\Omega_d$ is maximised. The natural zero is not an additional input. It is the unique structural feature of $p$ determined entirely by the recurrence constants.
\end{definition}

\subsection{Operation A: the zero-crossing forces $c_1$}

\begin{theorem}[Zero-crossing condition]
The natural zero condition $p(d_0) = 0$ forces $c_1 := q(d_0) = \frac{1}{2}\ln\pi$.
\end{theorem}

\begin{proof}
$p(d_0) = 0 \Leftrightarrow q(d_0) = \frac{1}{2}\ln\pi$. The right-hand side is the constant part of $p$, which by Theorem~3.1 equals $\frac{1}{2}\ln\pi$ with no free coefficient.
\end{proof}

\begin{theorem}[Uniqueness of Operation A]
Among all scalars derivable from the natural zero $d_0$, the unique scalar satisfying both (a) Range condition: $c>0$, so that $p(d) = c$ has a solution $d^*(c) > d_0$ in the cascade's descending regime, and (b) Primitivity condition: $c$ is a recurrence primitive, i.e., a $d$-independent constant appearing in the slicing recurrence with coefficient forced by orthogonality (Corollary~5.2), is $c_1 = q(d_0) = \frac{1}{2}\ln\pi$.
\end{theorem}

\begin{proof}
We show that only $q(d_0)$ passes both filters across all four information classes available from the natural zero.

Class 1: Values of $p$ and its decomposition at $d_0$. The zero gives three scalars: $p(d_0) = 0$, the $d$-dependent part $q(d_0) = \frac{1}{2}\ln\pi$, and the constant part $-\frac{1}{2}\ln\pi$. (i) $p(d_0) = 0$ fails (a): the equation $p(d) = 0$ has its unique solution at $d_0$ itself. (ii) $q(d_0) = \frac{1}{2}\ln\pi$ satisfies (a) since $\frac{1}{2}\ln\pi > 0$. It satisfies (b) since $\frac{1}{2}\ln\pi$ is the log-space image of $\sqrt{\pi}$, the recurrence's unique primitive. This is $c_1$. (iii) The constant part $-\frac{1}{2}\ln\pi < 0$ fails (a).

Class 2: Higher derivatives $p^{(n)}(d_0)$, $n\geq 1$. Using $\psi^{(n)}(x)\sim(n-1)!x^{-n}$ with $d_0 = 2\pi$: $p^{(n)}(d_0)\approx(n-1)!/[2(2\pi)^n]$. All derivatives fail (b): they are not recurrence primitives.

Class 3: The dimension $d_0$ itself. Setting $p(d) = d_0$ equates a dimensionless rate to a dimension count---a type mismatch.

Class 4: Analytic combinations. Any combination either reduces to $q(d_0)$ (Class 1), or introduces $\pi^{-n}$, $0$, or $d_0 = 2\pi$ as factors, none of which are recurrence primitives.

Across all four classes, $c_1 = q(d_0) = \frac{1}{2}\ln\pi$ is the unique scalar satisfying both (a) and (b).
\end{proof}

\subsection{Operation B: exp-conjugacy forces $c_2$}

\begin{theorem}[Exp-conjugacy]
The recurrence's linear form and log form represent the same underlying constant $\sqrt{\pi}$ in two different algebraic spaces. Their relationship forces $c_2 := e^{c_1} = \sqrt{\pi}$. $c_2$ is the linear-space image of $c_1$, derived from $c_1$ by the same exponential map that relates the two canonical forms of the recurrence.
\end{theorem}

\begin{proof}
The log-form constant is $c_1 = \frac{1}{2}\ln\pi$. The linear-form constant is $\sqrt{\pi}$. These represent the same mathematical object in two canonical representations: $e^{c_1} = e^{\frac{1}{2}\ln\pi} = \pi^{1/2} = \sqrt{\pi} = c_2$.
\end{proof}

\subsection{The orbit terminates}

\begin{theorem}[Orbit and termination]
Starting from $\{0\}$, define: Operation A (zero-crossing): $0\mapsto c_1 := q(d_0) = \frac{1}{2}\ln\pi$. Operation B (exp-conjugacy): $c_1\mapsto c_2 := e^{c_1} = \sqrt{\pi}$. The orbit of $\{0\}$ under $\{A,B\}$ is exactly $\{0, c_1, c_2\}$. No further application of $A$ or $B$ produces a new recurrence primitive.
\end{theorem}

\begin{proof}
We check all further applications: (i) $B$ applied to $0$: $e^0 = 1$. Rational; generates no threshold. (ii) $B$ applied to $c_2$: $e^{c_2} = e^{\sqrt{\pi}}\approx 5.83$. Does not appear anywhere in the recurrence derivation. (iii) $A$ applied to $c_1$ or $c_2$: Operation A acts on the unique natural zero of $p$. A strictly increasing function has at most one zero, and $p$ has exactly one (at $d_0$). There is no second zero at which an A-type operation could be applied. Therefore the orbit is $\{0, c_1, c_2\}$, and closure holds.
\end{proof}

\begin{corollary}[Canonical primitive set]
The canonical primitive set is $\mathcal{P} = \{0, c_1, c_2\} = \{0, \frac{1}{2}\ln\pi, \sqrt{\pi}\}$, derived from $\{0\}$ by operations $A$ and $B$, both internal to the recurrence.
\end{corollary}

\subsection{The two thresholds}

\begin{theorem}[First threshold: internal balance]
The condition $p(d) = c_1 = \frac{1}{2}\ln\pi$ selects $d_1 = 2\pi^2-1\approx 18.74$, integer $d_1 = 19$. At $d_1$, the volume ratio is $V_{d+1}/V_d = 1/\sqrt{\pi}$: each step loses exactly one factor of the recurrence's slicing constant.
\end{theorem}

\begin{proof}
\textit{Threshold selection.} Setting $p(d_1) = c_1$ and using $\psi(x)\approx\ln x$:
\[
\tfrac{1}{2}\ln\!\left(\frac{d_1+1}{2}\right) - \tfrac{1}{2}\ln\pi = \tfrac{1}{2}\ln\pi,
\]
giving $(d_1+1)/2 = \pi^2$, hence $d_1 = 2\pi^2-1$. Discrete verification: $p(19) = 0.5535 < c_1 = 0.5724 < p(20) = 0.6013$.

\textit{Volume ratio.} In the Stirling approximation,
\[
\ln\!\left(\frac{V_{d+1}}{V_d}\right) = \tfrac{1}{2}\ln\pi + \ln R(d) \approx \tfrac{1}{2}\ln\pi - \tfrac{1}{2}\ln\!\left(\frac{d+1}{2}\right),
\]
where the second step uses $\ln\Gamma((d+2)/2) - \ln\Gamma((d+1)/2)\approx\tfrac{1}{2}\ln((d+1)/2)$. The threshold condition fixes $\tfrac{1}{2}\ln((d_1+1)/2) = \ln\pi$. Substituting:
\[
\ln\!\left(\frac{V_{d_1+1}}{V_{d_1}}\right) \approx \tfrac{1}{2}\ln\pi - \ln\pi = -\tfrac{1}{2}\ln\pi,
\]
giving $V_{d_1+1}/V_{d_1} = 1/\sqrt{\pi}$.
\end{proof}

\begin{theorem}[Second threshold: cross-section cost]
The condition $p(d) = c_2 = \sqrt{\pi}$ selects $d_2 = 2\pi e^{2\sqrt{\pi}}-1\approx 216.6$, integer $d_2 = 217$.
\end{theorem}

\begin{proof}
Setting $p(d_2) = \sqrt{\pi}$ and using $\psi(x)\approx\ln x$: $(d_2+1)/2 = e^{\ln\pi + 2\sqrt{\pi}} = \pi e^{2\sqrt{\pi}}$, giving $d_2 = 2\pi e^{2\sqrt{\pi}}-1$. Discrete verification: $p(217) = 1.77101 < \sqrt{\pi} = 1.77245 < p(218) = 1.77331$.
\end{proof}

\subsection{Uniqueness of the two thresholds}

\begin{theorem}[Uniqueness]
No element of $\mathcal{P}$ other than $c_1$ and $c_2$ generates a distinct positive threshold. The cascade functional $p(d) = c$ admits exactly two solutions in the positive reals generated by $\mathcal{P}\setminus\{0\}$.
\end{theorem}

\begin{proof}
$\mathcal{P}\setminus\{0\} = \{c_1, c_2\}$ has exactly two elements by Theorem~6.5. Each is positive. Since $p$ is strictly monotone, each positive value generates exactly one threshold dimension.
\end{proof}

\section{The Tower of Distinguished Dimensions}

\subsection{The volume maximum and the sphere-area maximum}

Two distinct maxima arise from the cascade's geometry and should not be conflated.

The ball volume $V_d$ is maximised at $d=5$: $V_5\approx 5.26$ is the largest unit ball in any dimension. This is not generated by the threshold machinery of Section~6; no threshold condition $p(d) = c$ selects $d=5$. It is the argmax of $V_d$, a separate structural feature of the Gamma function.

The sphere area $\Omega_d$ is maximised where $p(d_0) = 0$, giving $d_0 = 2\pi\approx 6.28$, integer $d_0 = 7$. This is the natural zero of the recurrence---the unique structural feature of $p$ that seeds the orbit construction of Section~6 and generates the entire threshold pair $(d_1, d_2)$. The sphere-area maximum, not the volume maximum, is the cascade's primary distinguished point.

Both are recorded in Section~7.2 without physical interpretation. The separation between them---$V_d$ maximised at $d=5$, $\Omega_d$ maximised at $d=7$---is a consequence of the identity $\Omega_{d-1}/V_d = d$ (Theorem~4.7): boundary and volume measures peak at different dimensions precisely because they are related by a factor of $d$, not a constant.

\subsection{The tower}

\begin{center}
\begin{tabular}{cccl}
\toprule
Integer & Stirling & $\Omega_d$ & Meaning\\
\midrule
$d=5$ & $\operatorname{argmax}\,V_d$ & $\approx 26.3$ & Volume maximum\\
$d_0 = 7$ & $\approx 2\pi$ & $\approx 33$ & Sphere-area maximum (natural zero of $p$)\\
$d_1 = 19$ & $\approx 2\pi^2-1$ & $\approx 0.52$ & First threshold ($c_1 = \tfrac{1}{2}\ln\pi$)\\
$d_2 = 217$ & $\approx 2\pi e^{2\sqrt{\pi}}-1$ & $\approx 10^{-120}$ & Second threshold ($c_2 = \sqrt{\pi}$)\\
\bottomrule
\end{tabular}
\end{center}

\subsection{Three scales from one constant}

\begin{center}
\begin{tabular}{lccc}
\toprule
Scale & $\Omega_d$ & $\log_{10}\Omega$ & Ratio to $\Omega_7$\\
\midrule
$\Omega_7$ (sphere-area max) & 33.07 & 1.52 & 1\\
$\Omega_{19}$ (first threshold) & 0.516 & $-0.29$ & $1/64$\\
$\Omega_{217}$ (second threshold) & $2.3\times 10^{-120}$ & $-119.6$ & $\sim 10^{-121}$\\
\bottomrule
\end{tabular}
\end{center}

\begin{remark}[Scale variance from nothing]
The input to this construction is the concept of orthogonal dimensions. The output is three geometric scales spanning $10^{121}$. The mechanism is: natural zero of $p\to c_1\to c_2\to$ two threshold dimensions $\to$ two exponentially separated sphere areas. No mechanism for producing large hierarchies is assumed; the hierarchy is a property of the Gamma function evaluated at the recurrence's distinguished points.
\end{remark}

\section{The Inter-Threshold Identity}

\begin{theorem}[Threshold ratio]
$d_1/d_2 = e^{2(c_1-c_2)}$.
\end{theorem}

\begin{proof}
The cascade condition $p(d) = c$ gives $d = 2\pi e^{2c}-1\approx 2\pi e^{2c}$ in the Stirling regime. Therefore $d_1/d_2 = e^{2c_1}/e^{2c_2} = e^{2(c_1-c_2)}$. Substituting: $d_1/d_2 = \sqrt{\pi}\cdot e^{-2\sqrt{\pi}}$.
\end{proof}

\begin{corollary}[Inter-threshold fraction]
$(d_2-d_1)/d_2 = 1 - e^{-(2\sqrt{\pi}-\ln\pi)}$, a function of $\pi$ alone. Numerically: $198/217 = 0.91244$.
\end{corollary}

\begin{corollary}[Invariant as primitive function]
The cascade invariant is determined entirely by the primitive pair $(c_1, c_2)$: $I = \Omega_{d^*(c_1)}\times\Omega_{d^*(c_2)}\times(1-e^{2(c_1-c_2)})$. No quantity beyond the orbit output $\{c_1, c_2\}$ enters its definition.
\end{corollary}

\section{The Cascade Invariant}

\begin{definition}[Cascade invariant]
$I := \Omega_{d_1}\times\Omega_{d_2}\times(d_2-d_1)/d_2$.
\end{definition}

\textit{Sphere areas rather than volumes.} The thresholds $d_1, d_2$ are conditions on $p(d) = -\partial_d\ln\Omega_d$, which is defined as the log-derivative of sphere areas. Sphere areas are the primary objects of the recurrence; volumes are secondary. By Theorem~4.7, $\Omega_{d-1}/V_d = d$: at high $d$ the boundary dominates the volume, so sphere areas are the natural measure of the cascade's content at each layer.

\textit{Product rather than another symmetric function.} The slicing recurrence is multiplicative. The natural operation on two outputs of a multiplicative recurrence is their product.

\textit{Inter-threshold fraction.} The factor $(d_2-d_1)/d_2$ is the fraction of the cascade's total span occupied by the active inter-threshold region $[d_1, d_2]$. By Corollary~8.2, this fraction equals $1-e^{2(c_1-c_2)}$, determined entirely by the primitive pair.

\begin{theorem}[Uniqueness of the cascade invariant]
Under one geometric principle:
\begin{itemize}
\item[(A2)] \textit{Completeness of the primitive pair:} The invariant uses all information in the primitive pair $(c_1, c_2)$, including the inter-threshold fraction determined by Corollary~8.2,
\end{itemize}
the unique scalar associated to the primitive pair $(c_1, c_2)$ by the cascade's geometry is
\[
I = \Omega_{d_1}\times\Omega_{d_2}\times(d_2-d_1)/d_2.
\]
\end{theorem}

\begin{proof}
Five constraints, the first four from theorems already established, the fifth from (A2).

(1) \textit{Sphere areas, exponent one.} The slicing recurrence is a first-order multiplicative map outputting one copy of $\Omega_d$ at each step (Remark~4.11). Since $V_d = \Omega_{d-1}/d$ is a derived quantity, the only geometric measures available at each distinguished dimension are the sphere areas $\Omega_{d_1}$ and $\Omega_{d_2}$. Higher powers $\Omega^k$ ($k\neq 1$) would require raising to a power---an operation not present in the recurrence.

(2) \textit{Multiplicative combination.} The slicing recurrence is multiplicative. The natural operation combining outputs of a multiplicative system is their product.

(3) \textit{Two factors.} The orbit of the natural zero under operations $A$ and $B$ generates exactly the primitive pair $(c_1, c_2)$ and terminates (Theorem~6.5). Two distinguished dimensions, so two $\Omega$ factors.

(4) \textit{Product, not quotient.} The cascade originates at $d=\infty$ where $V_\infty = \Omega_\infty = 0$. As $d_2\to\infty$, the product $\Omega_{d_1}\times\Omega_{d_2}\to 0$. The quotient $\Omega_{d_1}/\Omega_{d_2}\to\infty$. Only the product has the correct limiting behaviour.

(5) \textit{Inter-threshold fraction (from A2).} The quantity $(d_2-d_1)/d_2 = 1-e^{2(c_1-c_2)}$ is determined by the primitives (Corollary~8.2). It encodes the geometric relationship between the two thresholds. Assumption (A2) requires that the invariant use all information in the primitive pair; including the inter-threshold fraction makes the invariant a property of the pair, not just the product of two separate quantities.

The unique scalar surviving all five constraints is $I = \Omega^{+1}_{d_1}\times\Omega^{+1}_{d_2}\times(d_2-d_1)/d_2$.
\end{proof}

\begin{theorem}[Evaluation]
$I = 1.0996\times 10^{-120}$.
\end{theorem}

\begin{proof}
Direct computation: $\Omega_{19} = 0.51614$, $\Omega_{217} = 2.33490\times 10^{-120}$, $198/217 = 0.91244$. Product: $1.0996\times 10^{-120}$.
\end{proof}

\subsection{Two geometric scales from two thresholds}

\begin{observation}
The sphere area at $d_1 = 19$ is $\Omega_{19} = 0.51614$---an $O(1)$ number in cascade units. Define: $G_1 := \sqrt{\Omega_{19}} = 0.71843$, $M_1 := (1/\Omega_{19})^{1/4} = 1.17980$. These are functions of $\pi$ alone, determined entirely by the first threshold $c_1 = \frac{1}{2}\ln\pi$.
\end{observation}

\begin{observation}
The sphere area at $d_2 = 217$, weighted by the inter-threshold fraction, is: $\Lambda_2 := \Omega_{217}\times 198/217 = 2.1310\times 10^{-120}$. This is a function of $\pi$ alone, determined by $c_2 = \sqrt{\pi}$ and the inter-threshold identity (Theorem~8.1).
\end{observation}

\begin{observation}
$I = \Omega_{19}\times\Lambda_2 = G_1^2\times\Lambda_2 = \Lambda_2/M_1^4$. Each factor is determined by its own threshold, independently.
\end{observation}

\begin{remark}[Numerical coincidences]
The cascade's two geometric scales bear a striking numerical resemblance to two fundamental constants of physics. The observed ratio of vacuum energy to the fourth power of the Planck mass is $\rho_\Lambda/M^4_{\mathrm{Pl}}\approx 1.1\times 10^{-120}$, which matches $I = \Lambda_2/M_1^4$ to 0.04\%. The physical and geometric framework within which this identification may be derived lies outside the scope of this paper.
\end{remark}

\section{The 120 Orders of Magnitude}

\begin{theorem}[The hierarchy]
$\log_{10}(1/\Omega_{d_2})\approx\pi e^{2\sqrt{\pi}}(2\sqrt{\pi}-1)/\ln 10\approx 120.3$.
\end{theorem}

\begin{proof}
Stirling at $d_2 = 2\pi e^{2\sqrt{\pi}}$: $\ln\Omega_{d_2}\approx -\pi e^{2\sqrt{\pi}}(2\sqrt{\pi}-1)\approx -277$ nats $\approx -120.3$ decades. The mechanism is the Gamma function's superexponential growth: as $d$ increases from 19 to 217, the denominator $\Gamma(d/2)$ in $\Omega_d$ grows from $\Gamma(9.5)\approx 10^5$ to $\Gamma(108.5)$, a number with over 170 digits.
\end{proof}

\begin{remark}[Anatomy of the hierarchy]
$\sqrt{\pi}\approx 1.77$ (orthogonal compression constant); $e^{2\sqrt{\pi}}\approx 34.6$ (exponential amplification); $2\pi e^{2\sqrt{\pi}}\approx 218$ (the cascade dimension); $(2\sqrt{\pi}-1)\approx 2.54$ (per-dimension suppression). Product: $\approx 277$ nats $\approx 120$ decades.
\end{remark}

\section{Comparison with a Physical Constant}

The cascade invariant $I = 1.0996\times 10^{-120}$ is a pure number from the Gamma function. It was not constructed to match any physical quantity.

The observed ratio of the vacuum energy density to the fourth power of the Planck mass, measured by cosmological observations, is $\rho_\Lambda/M^4_{\mathrm{Pl}}\approx 1.1\times 10^{-120}$. This number is determined by general relativity, quantum field theory, and astronomical measurements---none of which enter the present paper. The agreement with $I$ is 0.04\% under the standard reduced-Planck-mass convention.

We record this agreement without explaining it.

\section{Summary of the Derivation}

\begin{center}
\begin{tabular}{cll}
\toprule
Step & Statement & Source\\
\midrule
1 & $V_\infty = \Omega_\infty = 0$; rates are primary & Definition\\
2 & Orthogonal $\Rightarrow$ quarter turn $\Rightarrow B(1/2,\cdot)$ & Def.\ of dimension\\
3 & $\Gamma(1/2) = \sqrt{\pi}$: unique constant, coeff.\ 1; irreducible & Thms 3.1, 5.1\\
3A & Each step: $R_{\mathrm{eff}} = 1/\sqrt{d+3}$; $\Omega/V = d$; $V_d = \Omega_{d-1}/d$ & Thms 4.2--4.12\\
3B & $R_{\mathrm{eff}}/\sigma\to 1$; $\operatorname{erf}(1/\sqrt{2})$ is second cascade constant & Thm 4.4, Cor 4.5\\
3C & $\Omega_{d-1}$ is the unique independent cascade quantity & Cor 4.12\\
4 & $p(d)$ has unique natural zero $d_0$ (strict monotonicity) & $p'>0$\\
5 & Orbit seed $\{0\}$: value of $p$ at unique structural zero $d_0$ & Rem 2.3\\
6 & Zero-crossing forces $c_1 = \tfrac{1}{2}\ln\pi$; unique valid reading & Thms 6.2, 6.3\\
7 & Exp-conjugacy forces $c_2 = \sqrt{\pi}$; orbit terminates & Thms 6.4, 6.5\\
8 & Exactly two positive thresholds: $d_1 = 19$, $d_2 = 217$ & Thms 6.7--6.9\\
9 & $d_1/d_2 = e^{2(c_1-c_2)}$; inter-threshold fraction determined & Thm 8.1\\
10 & $I = \Omega(c_1)\times\Omega(c_2)\times(1-e^{2(c_1-c_2)})$; zero free params & Cor 8.3\\
11 & $I = 1.0996\times 10^{-120}$ & Direct computation\\
12 & $I = \Lambda_2/M_1^4$; each factor from its own threshold & Obs 9.4--9.6\\
13 & $120\approx\pi e^{2\sqrt{\pi}}(2\sqrt{\pi}-1)/\ln 10$ & Thm 10.1\\
14 & Combination is the unique invariant under (A2) & Thm 9.2, Rem 4.11\\
\bottomrule
\end{tabular}
\end{center}

The derivation from orthogonality to $I$ is complete under one geometric principle (A2: completeness of the primitive pair) stated in Theorem~9.2. The constraint that sphere areas enter with exponent 1 follows from the recurrence's first-order multiplicative structure (Remark~4.11); it is not an independent assumption. The threshold pair $\{c_1, c_2\}$ is the orbit of the natural zero $\{0\}$ under two operations internal to the recurrence; it requires no external definition. The invariant $I$ has zero free parameters beyond the orbit output (Corollary~8.3), and its uniqueness is established under (A2).

\section{What This Paper Does Not Do}

\begin{itemize}
\item It does not explain the agreement between $I$ and $\rho_\Lambda/M^4_{\mathrm{Pl}}$. It records the agreement as an empirical observation.
\item The uniqueness of the invariant (Theorem~9.2) is established under one geometric principle (A2: completeness of the primitive pair). The constraint that sphere areas enter with exponent 1 follows from the recurrence's multiplicative structure (Remark~4.11) and is not an independent assumption. A proof that (A2) is the only consistent reading of primitive completeness would make the uniqueness unconditional.
\item It does not derive general relativity, quantum mechanics, the Standard Model, or any physical law.
\item It does not explain why we observe four dimensions.
\item It does not provide a physical mechanism for the scale variance.
\item It uses no result from physics. The paper is a theorem about the Gamma function.
\item The layered boundary dominance (Corollary~4.10) is a mathematical statement about Gamma function ratios. This paper assigns no physical significance to it.
\item Corollary~4.12 establishes that sphere areas are the only independent cascade quantities, and that any physical identification mapping cascade content to energy must do so linearly with a universal constant. The derivation of this consequence lies outside the scope of this paper.
\end{itemize}

\section{Robustness Test}

To test whether the cascade invariant is an isolated numerical coincidence or a unique output of the tower structure, we systematically evaluated 249 distinct combinations of the cascade's outputs ($d_0, d_1, d_2, \Omega_7, \Omega_{19}, \Omega_{217}$). Result: exactly one combination falls within 1\% of the observed value: $\Omega_{19}\times\Omega_{217}\times 198/217 = 1.0996\times 10^{-120}$ ($-0.04\%$). The next closest is $\Omega_{19}\times\Omega_{217} = 1.2051\times 10^{-120}$ ($+9.6\%$). Shifting either threshold by one integer gives deviations of 70--900\%. The match is sharply localised at $(d_1, d_2) = (19, 217)$.

\begin{center}
\begin{tabular}{llcl}
\toprule
Primitive set & Thresholds & $\log_{10} I$ & Deviation\\
\midrule
$\{c_1, c_2\}$ (this paper) & 19, 217 & $-119.96$ & 0.04\%\\
$\{0, c_2\}$ & 7, 217 & $-118.12$ & Factor 68\\
$\{c_1, 2c_1\}$ & 19, 61 & $-17.3$ & $10^3$ OOM\\
$\{2c_1, c_2\}$ & 61, 217 & $-136.6$ & 17 OOM low\\
$\{c_1/2, c_2\}$ & 36, 217 & $-126.4$ & 6 OOM\\
$\{c_1, \pi\}$ & 19, $>$3000 & $<-1000$ & Irrelevant\\
\bottomrule
\end{tabular}
\end{center}

No alternative primitive set drawn from operations present in the recurrence comes within an order of magnitude of $10^{-120}$. The 0.04\% agreement is unique to the orbit-generated pair.

\appendix

\section{Numerical Verification}

Constants: $\sqrt{\pi} = 1.7724538509\ldots$, $\frac{1}{2}\ln\pi = 0.5723649429\ldots$

Continuous distinguished dimensions: $d_0 = 2\pi = 6.283$, $d_1 = 2\pi^2-1 = 18.739$, $d_2 = 2\pi e^{2\sqrt{\pi}}-1 = 216.63$.

Discrete first threshold: $p(19) = 0.5535 < c_1 = 0.5724 < p(20) = 0.6013$. Confirmed.

Discrete second threshold: $\Delta(d) := p(d) - \sqrt{\pi}$: $\Delta(217) = -0.00144 < 0 < \Delta(218) = +0.00086$. That is, $p(217) = 1.77101 < \sqrt{\pi} = 1.77245 < p(218) = 1.77331$. Confirmed.

Scale coincidence: $R_{\mathrm{eff}}(d)/\sigma(d) = \sqrt{d/(d+3)}$. At $d=217$: $\sqrt{217/220} = 0.99318$. At $d=\infty$: 1 exactly. The band area fraction at $d=217$: $F(217) = \operatorname{erf}(\sqrt{217/440}) = 0.6821$; at $d=\infty$: $\operatorname{erf}(1/\sqrt{2}) = 0.68269$.

Sphere areas: $\Omega_7 = 32.47$, $\Omega_{19} = 0.51614$, $\Omega_{217} = 2.33490\times 10^{-120}$. Cascade invariant: $\Omega_{19}\times\Omega_{217} = 1.20513\times 10^{-120}$; $198/217 = 0.91244$; $I = 1.09961\times 10^{-120}$. Observed: $\rho_\Lambda/M^4_{\mathrm{Pl}} = 1.1\times 10^{-120}$. Deviation: 0.04\%.

Hierarchy: $\pi e^{2\sqrt{\pi}}(2\sqrt{\pi}-1)/\ln 10 = 120.27$.

Independent geometric scales: $G_1 = \sqrt{\Omega_{19}} = 0.71843$; $M_1 = (1/\Omega_{19})^{1/4} = 1.17980$; $\Lambda_2 = \Omega_{217}\times 198/217 = 2.1310\times 10^{-120}$; $\Lambda_2\times\Omega_{19} = 1.0997\times 10^{-120} = I$. $\checkmark$

Compactification radii: $R_{\mathrm{eff}}(4) = 1/\sqrt{7} = 0.378$, $R_{\mathrm{eff}}(12) = 1/\sqrt{15} = 0.258$, $R_{\mathrm{eff}}(19) = 1/\sqrt{22} = 0.213$, $R_{\mathrm{eff}}(217) = 1/\sqrt{220} = 0.067$.

\section{Sign Pattern of the Odd Logarithmic Derivatives of $\xi$}

The following theorem is a consequence of the cascade's threshold structure. It does not assume the Riemann Hypothesis and requires no analytic continuation beyond the real axis.

\begin{theorem}[Alternating sign pattern of odd derivatives]
Let $\xi(s)$ denote the completed Riemann zeta function and let $\kappa_{2k+1}(\sigma) := \partial^{2k+1}_\sigma\log\xi(\sigma)$ for $\sigma\in(\frac{1}{2},1)$ and $k\geq 0$. Then $\operatorname{sgn}(\kappa_{2k+1}(\sigma)) = (-1)^k$ for all $k = 0, 1, \ldots, 43$ and all $\sigma\in(\frac{1}{2},1)$, as a consequence of the single inequality $\gamma_1 > (2+\sqrt{3})/2\approx 1.87$, where $\gamma_1\approx 14.13$ is the imaginary part of the first non-trivial zero. The result assumes no zeros exist outside the standard zero-free region; it does not assume the Riemann Hypothesis.
\end{theorem}

\begin{proof}
From the Hadamard product for $\xi$, each non-trivial zero pair $\rho = \beta+i\gamma$, $\bar{\rho} = \beta-i\gamma$ contributes to $\kappa_{2k+1}(\sigma)$ a term proportional to $2\,\mathrm{Re}[(\sigma-\rho)^{-(2k+1)}]$. Writing $u = \sigma-\beta$, $\gamma = r\sin\theta$ with $\theta = \arctan(\gamma/u)$: the sign is $\operatorname{sgn}(\cos((2k+1)\theta))$. Within the standard zero-free region all zeros satisfy $|\beta - \sigma| \leq \frac{1}{2}$ for $\sigma\in(\frac{1}{2},1)$, so $\theta\in[\arctan(\gamma_1),\pi/2)\subset(1.499, 1.571)$. The condition that no zero of $\cos((2k+1)\theta)$ falls in this interval holds for $k\leq 43$ because the interval of length $0.046(2k+1)$ contains no odd integer for these values. This fails at $k=44$, where 87 first enters the problematic range. Since $\cos((2k+1)\theta) > 0$ for all $\theta$ in the zero's angular range, and the archimedean and pole terms also have sign $(-1)^k$, the total $\kappa_{2k+1}(\sigma)$ has sign $(-1)^k$ for $k\leq 43$.
\end{proof}

\begin{remark}
The bound $k\leq 43$ is sharp: the proof fails for $k=44$. The threshold integer 43 is the Heegner number following 19 in the sequence $\{7, 11, 19, 43, 67, 163\}$. Note that 7 and 19 are cascade threshold dimensions and 43 is the cumulant sign bound; 11 plays no cascade role, so the alignment is partial and noted as a coincidence rather than a structural connection.
\end{remark}

\begin{thebibliography}{9}

\bibitem{cascadeII} [Author], \textit{Quantum Mechanics from the Cascade: Effective Theory of a 4-Dimensional Observer in the Sphere-Area Geometry}, March 2026.

\bibitem{cascadeIII} [Author], \textit{General Relativity from the Cascade: Lovelock Uniqueness and the Cosmological Constant of a 4-Dimensional Observer}, March 2026.

\bibitem{cascadeIVa} [Author], \textit{The Standard Model from the Cascade: Part IVa---Gauge Group, Symmetry Breaking, and Three Generations}, March 2026.

\bibitem{cascadeIVb} [Author], \textit{The Standard Model from the Cascade: Part IVb---Masses, Couplings, and Precision Predictions}, March 2026.

\bibitem{ledoux} M. Ledoux, \textit{The Concentration of Measure Phenomenon}, AMS Mathematical Surveys and Monographs 89, 2001.

\bibitem{ball} K. Ball, ``An elementary introduction to modern convex geometry,'' in \textit{Flavors of Geometry}, MSRI Publications 31, Cambridge, 1997.

\bibitem{weinberg} S. Weinberg, ``The cosmological constant problem,'' \textit{Reviews of Modern Physics} \textbf{61}, 1--23 (1989).

\bibitem{padmanabhan} T. Padmanabhan, ``Some aspects of the cosmological constant problem,'' \textit{Physics Reports} \textbf{380}, 235--320 (2003).

\bibitem{desi} DESI Collaboration, ``DESI 2024 VI: Constraints on models of dark energy,'' arXiv:2404.03002 (2024).

\end{thebibliography}

\end{document}
